Considera la funció
\[
g(x)=\begin{cases}
  \dfrac{x^2-1}{x+1} & \si{x<2},\\[8pt]
  3x+k & \si{x\ge 2},
\end{cases}
\]
on $k$ és un nombre real.

\begin{apartats}

\apartat[1]{0,75}
Calcula els límits següents:
\begin{graella}{2}
  \sa \lim_{x\to-1}g(x) & \sa \lim_{x\to0}g(x)
\end{graella}

\begin{solucio}
Tots dos punts són a la branca $x<2$.\\
En $x=-1$ hi ha una indeterminació $0/0$:
$\dfrac{x^2-1}{x+1}=\dfrac{(x+1)(x-1)}{x+1}=x-1\to-2$.
El límit val $-2$, tot i que $g(-1)$ no existeix.\\
En $x=0$ se substitueix directament: $\dfrac{-1}{1}=-1$.
\end{solucio}

\apartat[1,5]{1,25}
Calcula, en funció de $k$, els límits laterals de $g$ en $x=2$. Per a quin valor de $k$
existeix $\lim_{x\to2}g(x)$? Quant val aquest límit?

\begin{solucio}
$\lim_{x\to2^-}g(x)=\dfrac{4-1}{2+1}=1$ \quad i \quad $\lim_{x\to2^+}g(x)=6+k$.\\
El límit existeix si i només si els dos laterals coincideixen: $6+k=1$, és a dir
$\boxed{k=-5}$. Aleshores $\lim_{x\to2}g(x)=1$.
\end{solucio}

\begin{nomesllarg}
\apartat{0,5}
Per a $k=-5$, calcula els límits següents:
\begin{graella}{2}
  \sa \lim_{x\to-\infty}g(x) & \sa \lim_{x\to+\infty}g(x)
\end{graella}

\begin{solucio}
Quan $x\to-\infty$ actua la primera branca, que per a $x\ne-1$ és $x-1\to-\infty$.\\
Quan $x\to+\infty$ actua la segona: $3x-5\to+\infty$.
\end{solucio}
\end{nomesllarg}

\end{apartats}
